% Two-round Bayesian arm choice. All numbers are exact teaching values.
\begin{tikzpicture}[x=1mm,y=1mm,
  every node/.style={font=\figurefont\fontsize{8}{10}\selectfont,text=FigureInk,inner sep=1mm},
  action/.style={draw=FigureBlue,rounded corners=1.8pt,line width=.7pt,align=center,fill=white,text width=25mm,minimum height=13mm},
  observation/.style={draw=FigureStroke,rounded corners=1.8pt,line width=.6pt,align=center,fill=white,text width=19mm,minimum height=10mm},
  posterior/.style={draw=FigureGreen,rounded corners=1.8pt,line width=.7pt,align=center,fill=FigureGreenFill,text width=31mm,minimum height=10mm},
  nextaction/.style={draw=FigureBlue,rounded corners=1.8pt,line width=.7pt,align=center,fill=white,text width=35mm,minimum height=10mm},
  branch/.style={-{Stealth[length=1.8mm,width=1.2mm]},draw=FigureStroke,line width=.7pt},
  note/.style={font=\FigureNoteFont\fontsize{7.5}{9.5}\selectfont,text=FigureMuted}]
  \path[use as bounding box] (0,0) rectangle (169,83);
  \node at (16,79) {第1轮行动};
  \node at (55,79) {第1轮观测};
  \node at (93,79) {参数后验};
  \node at (142,79) {第2轮最优行动};
  \node[action] (known) at (16,62) {已知臂\\成功率\(0.6\)};
  \node[observation] (known-result) at (55,62) {任一结果};
  \node[posterior] (unchanged) at (93,62) {高参数\(0.9\)的概率\\\(0.5\)};
  \node[nextaction] (tie) at (142,62) {两臂仍并列\\成功率\(0.6\)};
  \draw[branch] (known)--(known-result);
  \draw[branch] (known-result)--(unchanged);
  \draw[branch] (unchanged)--(tie);
  \node[action] (unknown) at (16,25) {未知臂\\成功率\(0.6\)};
  \node[observation] (success) at (55,38) {成功，概率\(0.6\)};
  \node[observation] (failure) at (55,12) {失败，概率\(0.4\)};
  \node[posterior] (high) at (93,38) {高参数\(0.9\)的概率\\\(0.75\)};
  \node[posterior] (low) at (93,12) {高参数\(0.9\)的概率\\\(0.125\)};
  \node[nextaction] (choose-unknown) at (142,38) {选未知臂\\成功率\(0.75\)};
  \node[nextaction] (choose-known) at (142,12) {选已知臂\\成功率\(0.6\)};
  \draw[branch] (unknown)--(success);
  \draw[branch] (unknown)--(failure);
  \draw[branch] (success)--(high);
  \draw[branch] (failure)--(low);
  \draw[branch] (high)--(choose-unknown);
  \draw[branch] (low)--(choose-known);
  \node[note] at (142,1.5) {失败后未知臂成功率为\(0.375\)};
\end{tikzpicture}
